\documentclass{article}
\usepackage{graphicx} % Required for inserting images
\usepackage[english]{babel}
\usepackage{amsfonts}
\usepackage{amssymb}
\usepackage{amsmath}
\usepackage{a4wide}
\usepackage{xcolor}
\usepackage{enumitem}


\title{Stat 17 Week 4}
\author{Xiao Zhang}
\date{January 2026}

\begin{document}

\maketitle



\section*{Probability Topics — Part 2: Events \& Probability Rules}

\begin{itemize}
  \item \textbf{Events}
  \begin{itemize}
    \item An event is any collection of outcomes from the sample space
    \item Impossible events have probability 0
    \item Certain events have probability 1
    \item Events with very small probability (often $<0.05$) are considered unusual
  \end{itemize}

  \item \textbf{Basic Probability Rules}
  \begin{itemize}
    \item $0 \le P(A) \le 1$ for any event $A$
    \item The probabilities of all outcomes in the sample space sum to 1
  \end{itemize}

  \item \textbf{Complements}
  \begin{itemize}
    \item The complement of event $A$, denoted $A^c$, includes all outcomes not in $A$
    \item Complement rule:
    \[
      P(A^c) = 1 - P(A)
    \]
  \end{itemize}


  \item \textbf{Unions and Intersections}
  \begin{itemize}
    \item Intersection: $A \cap B$ (both $A$ and $B$ occur)
    \item Union: $A \cup B$ ($A$ or $B$ or both occur)
  \end{itemize}

  \item \textbf{Mutually Exclusive Events}
  \begin{itemize}
    \item Two events are mutually exclusive if they share no outcomes
    \item If $A$ and $B$ are mutually exclusive:
    \[
      P(A \cup B) = P(A) + P(B)
    \]
  \end{itemize}

  
  \item \textbf{Addition Rule (OR)}
  \begin{itemize}
      \item 
For any two events $A$ and $B$,
\[
P(A \cup B) = P(A) + P(B) - P(A \cap B).
\]

    \item If $A$ and $B$ are \textbf{mutually exclusive} (cannot occur together, i.e. $P(A \cap B) = 0$), then
\[
P(A \cup B) = P(A) + P(B).
\]
  \end{itemize}

  \item \textbf{Decomposition of an Event}
    Any event $A$ can be written as the union of two disjoint events:
    \[
    A = (A \cap B) \cup (A \cap B^c).
    \]
    
    Since $(A \cap B)$ and $(A \cap B^c)$ are mutually exclusive,
    \[
    P(A) = P(A \cap B) + P(A \cap B^c).
    \]



  \item \textbf{Conditional Probability}
  \begin{itemize}
    \item The probability of $A$ given $B$ is:
    \[
      P(A|B) = \frac{P(A \cap B)}{P(B)}
    \]
    \item The ``given'' event always appears in the denominator
  \end{itemize}

  \item \textbf{Multiplication Rule(AND)}
  \begin{itemize}
    \item  For any two events $A$ and $B$,
    \[P(A \cap B) = P(A \mid B)\,P(B).\]
    \itme If $A$ and $B$ are \textbf{independent}, then
    \[P(A \cap B) = P(A)\,P(B).\]
  \end{itemize}

  \item \textbf{Independent Events}
  \begin{itemize}
    \item Events $A$ and $B$ are independent if one does not affect the probability of the other
    \item Equivalent conditions:
    \begin{itemize}
      \item $P(A|B) = P(A)$
      \item $P(B|A) = P(B)$
      \item $P(A \cap B) = P(A)P(B)$
    \end{itemize}
  \end{itemize}
  


  \item \textbf{At Least One Events}
  \begin{itemize}
    \item These are often easiest to compute using complements:
    \[
      P(\textit{at least one}) = 1 - P(\textit{none})
    \]
  \end{itemize}



  \item \textbf{Law of Total Probability:}
  If $\{B_1, B_2, \dots, B_n\}$ is a partition of the sample space with $P(B_i)>0$, then
    \[
    P(A) = \sum_{i=1}^n P(A \mid B_i)\,P(B_i).
    \]

  \item \textbf{Tree Diagrams and Bayes' Rule}
  \begin{itemize}
    \item Tree diagrams help visualize conditional and joint probabilities
    \item Bayes' Rule allows probability updates given new information
  \end{itemize}
\end{itemize}



\end{document}
